\section{Literature Review}

\subsection{Surveying the Field}

The theory behind why models fail to generalise is well developed. Ben-David et al. (2010) proved mathematically that a model's ability to transfer from one population to another depends on how different the two populations are. When the difference is large enough, no algorithm  no matter how sophisticated  can make up for it. Gretton et al. (2012) built on this by creating the Maximum Mean Discrepancy (MMD) test, a statistical method for measuring exactly how different two datasets are. This matters because you need to detect a mismatch before you can do anything about it.
These theoretical limits show up in practice across many areas. In genomics, Whalen et al. (2022) mapped out five common mistakes that cause models to look better than they really are: using data that is not truly independent, allowing information to leak from test data into training, ignoring confounders, failing to account for class imbalance, and testing on distributions that do not match real-world deployment. They backed up each pitfall with interactive computational notebooks that let researchers reproduce the problem on their own data. In clinical biomarker research, Wörheide et al. (2021) made an important distinction between two types of data shift  covariate shift (where the input features change) and prior probability shift (where the disease prevalence changes)  and showed that mixing up the two leads to applying the wrong correction. They found that a common fix called importance weighting works for simple models but actually makes things worse for flexible models like gradient boosting. They also pointed to distributionally robust optimisation (DRO) as a promising alternative, referencing work by Duchi and Namkoong (2018) on KL-based approaches and Rahimian and Mehrotra (2019) on Wasserstein-based approaches, though importantly they recommended DRO without ever testing it themselves. In immunology, Springer and Doering (2021) showed that models predicting T-cell receptor binding performed well on targets they had seen during training but dropped to chance level on genuinely new targets  and that the overall performance metric hid this failure because it was dominated by a few well-represented targets. Chekroud et al. (2024) provided the most dramatic example: their schizophrenia treatment models went from looking highly accurate within each trial to being effectively useless across trials.
Several responses have been proposed. Ektefaie et al. (2024) created SPECTRA, a framework that tests models across the full range of overlap between training and test data, and found that all 19 deep learning models they tested got worse as the overlap decreased. Oloruntoba et al. (2023) showed that using more realistic data splits produces much lower performance estimates than the random splits most studies rely on. Pavlovic et al. (2024) argued that the real problem is models learning spurious correlations instead of genuine causal patterns, and showed that using causal graphs to guide feature selection produces more stable results  though their method requires domain experts to build the causal graph in advance. Arjovsky et al. (2019) proposed invariant risk minimisation, which looks for features that predict the outcome consistently across multiple training environments, but this requires those multiple environments to exist, which clinical datasets rarely provide. On the reporting side, Collins et al. (2024) developed the TRIPOD+AI checklist  a 27-item standard for how prediction model studies should be reported  but as a transparency tool rather than a quality tool, a model can tick every box and still fail at deployment. Other contributions from Lever et al. (2016), Okser et al. (2014), Runcie and Cheng (2019), Dincer et al. (2022), Evans et al. (2021), and Steyerberg et al. (2010) each show domain-specific versions of the same underlying problem: the way we currently evaluate models systematically makes them look better than they are.


\subsection{How the Field Frames the Problem}

The papers reviewed here approach generalisation failure from different angles, and each angle has a blind spot.
One set of papers sees the problem as one of good practice. Whalen et al. (2022) list avoidable pitfalls, Collins et al. (2024) specify what should be reported, Riley et al. (2021) work out sample size requirements, and Vayena et al. (2018) cover ethical considerations. The shared logic is simple: if researchers follow the right steps and report everything properly, models will generalise. This is helpful advice, but it cannot account for failures that happen even when the methodology is correct  which is exactly what Chekroud et al. (2024) show happening across five well-designed clinical trials.
A second set of papers accepts that the problem runs deeper  the data a model trains on is genuinely different from the data it will meet in practice. Ben-David et al. (2010) showed mathematically that there is a point where the two populations are so different that no fix will work. Wörheide et al. (2021) showed that common corrections like importance weighting depend on which model you are using  they help simple models but hurt complex ones. Springer and Doering (2021) showed that when you try to predict something the model has literally never seen before, it fails completely. Subbaswamy and Saria (2020) argue the only lasting solution is to find features that are causally linked to the outcome, not just statistically associated. The tension between Ben-David et al.'s mathematical pessimism and the more hopeful applied papers is real  and neither side has tested their position against the other using the same dataset.
A third group makes the most challenging argument: the problem is not the data or the methodology, but the way machine learning works at its core. Arjovsky et al. (2019) and Pavlovic et al. (2024) point out that standard ML simply learns whatever patterns reduce training error, including patterns that are accidents of the training population and will not appear anywhere else. Causal DAGs and invariant risk minimisation try to fix this by looking for features that predict the outcome consistently across different settings. Watson and Szathmary (2016) take this further, arguing through learning theory that the generalisation problem in ML mirrors how biological evolution transfers adaptations across environments. The practical obstacle is clear: IRM needs multiple separate training environments that clinical data almost never has, and building a causal graph requires specialist knowledge most ML teams do not possess.
Running underneath all of these is a conceptual point that Sullivan (2023) makes clearly: there is a difference between statistical generalisation (the model works on held-out data from the same population) and practical robustness (the model works when the population itself changes). Cross-validation tests the first. Clinical deployment requires the second. The field routinely treats evidence for one as proof of the other, and that is a mistake. Gretton et al. (2012), through the MMD test, offer one of the few existing tools for measuring how large this gap actually is.
The clinical papers  Chekroud et al. (2024) and Steyerberg et al. (2010)  show what happens at the point of deployment. Steyerberg et al. define what good performance looks like (discrimination, calibration, clinical usefulness). Chekroud et al. show how badly real models fall short. Collins et al. (2024) ensure that at least these failures are reported transparently. But transparency is where the current work stops. The literature is good at explaining failures after they happen. It has almost nothing to say about predicting them in advance.


\subsection{Deep Critical Analysis of Six Foundation Papers}

Chekroud et al. (2024), Science. Chosen for its recency, the authority of the journal, and the strength of its evidence  five international clinical trials accessed through the YODA Project. The performance collapse from 0.72 (within-trial) to 0.60 (cross-validated) to 0.54 (independent trial) holds across four different ways of defining treatment success and two different algorithms. That consistency makes it hard to dismiss. The main weakness is scope: schizophrenia treatment response may be an unusually difficult prediction problem, so the jump from "these models fail" to "clinical models generally fail" is bigger than the paper fully justifies. The design also cannot tell you whether the failure is caused by overfitting (learning noise) or by transportability failure (learning real but context-specific patterns)  and these two problems need completely different solutions. Most importantly for this portfolio, the authors never test any method designed to improve robustness. They show the problem clearly but leave the question of whether it can be solved completely open.
Whalen et al. (2022), Nature Reviews Genetics. Chosen for its clear organisation, practical demonstrations, and the reproducible notebooks that come with it. Their five-pitfall framework is directly useful  they show that dependent examples can inflate performance metrics by more than 0.5, and that leaky preprocessing can make a model appear accurate on data that contains no real signal at all. The limitation is that their framework assumes generalisation will follow if you fix these technical problems. But Chekroud et al.'s models appear to fail for reasons that go beyond methodology  the statistical patterns in the data simply do not transfer between trials. The paper's treatment of confounding is also limited: adjusting for confounders with PCA removes the visible effect but does not guarantee the model has learned relationships that will hold up in new settings.
Pavlovic et al. (2024), Nature Machine Intelligence. Chosen because it is the closest thing in the literature to a practical solution rather than just a diagnosis. Their causal graph approach identifies which features are genuinely linked to the outcome and which are confounded, and their simulations show that confounder shift roughly doubles error rates and that selection bias can make a model perform at chance on unbiased data. The key limitation is that every experiment uses a perfectly specified causal graph  one where all the cause-and-effect relationships are known in advance and correct. In real clinical settings, that knowledge does not exist. The critical unanswered question is what happens when the graph is only partly right: does a rough causal graph still beat a standard model, or does it actually make things worse? That question is the biggest barrier to anyone actually using this approach in practice.
Springer and Doering (2021), Briefings in Bioinformatics. Chosen because it tests what happens at the extreme edge of generalisation  predicting outcomes for biological targets that were completely absent from training. Their most revealing finding is the decoy experiment: published models that appeared to predict TCR-epitope binding actually performed just as well when the epitopes were replaced with random nonsense sequences. This proves those models were memorising receptor patterns rather than learning genuine interaction biology. Their own model can extrapolate, but only to epitopes within about one amino acid change of something in the training data  which is really interpolation at the boundary, not true generalisation to the unknown.
Collins et al. (2024), TRIPOD+AI, BMJ. Chosen for its authority as the field's reporting standard  built through a structured process with 200 international participants and registered with the EQUATOR Network. The 27-item checklist is a genuine step forward, embedding fairness requirements across twelve items and requiring open science practices. The fundamental limitation is that it is a transparency tool, not a quality tool. A model could satisfy every single TRIPOD+AI requirement and still perform at chance level on new patients. Item 26 asks researchers to discuss generalisation limitations, but sets no bar for what that discussion should contain  writing "this model may not generalise" counts the same as a detailed transportability analysis. If the standard cannot distinguish adequate from inadequate, it cannot protect against deployment failure.

\subsection{Synthesis}

Three things the literature agrees on: internal validation alone does not prove a model will generalise; dataset shift is present in virtually every real-world deployment scenario but the field has not settled on how to handle it; and spurious correlations  where the model learns patterns that are real in the training data but will not hold elsewhere  cannot be fixed with standard statistical techniques alone.
Two disagreements remain unresolved. The first is between researchers who see generalisation failure as a data quality problem (fix the methodology and the model will work) and those who see it as a structural problem (even perfect methodology will fail if the model learns the wrong patterns). Chekroud et al. use standard methods and watch them fail but never try causal approaches; Pavlovic et al. show causal methods help but never compare them against well-implemented standard corrections. The experiment that would settle this  testing both approaches on the same data  has not been done. The second disagreement is between the practical optimism of applied researchers who believe corrections can work and the theoretical caution of Ben-David et al., who show that beyond a certain level of distributional difference, nothing can help. The field does not yet have a way to tell which situation you are in for any given deployment.
Three gaps stand out. First, no method exists that can predict before deployment whether a model will fail in a specific new context  the literature can only explain failures after they have happened. Second, no paper treats post-deployment monitoring as a primary research problem: how do you detect that a deployed model is starting to degrade, and when should you intervene? Third, the most promising solution  Pavlovic et al.'s causal framework  requires specialist expertise that most clinical ML teams simply do not have.
These gaps point to a single missing piece. The field knows models fail, understands why, has theoretical fixes, and has reporting standards. What it lacks is a practical framework that can estimate the performance gap before deployment, identify which features are reliable and which are not, give clear guidance on what to do, and keep monitoring after deployment  all without needing a causal graph, multiple training environments, or a team of specialists. That missing piece  a gap in assembly, not in awareness  is what this project aims to provide.
